\section{Introduction}
\label{sec:intro}

Street-view perception models predict subjective attributes such as
safety, wealth, or liveliness at scale~\cite{salesses2013,naik2014,dubey2016},
yet remain limited as explanatory tools: they do not show which
localized visual changes would plausibly shift human judgement for a
specific scene. Existing explainability methods (saliency, SHAP,
concept probes) remain correlational, identifying features associated
with a score without validating that manipulating them changes human
judgement~\cite{adebayo2018,kim2018}. Urban-design research links cues
such as greenery, maintenance, and visibility to perceived
safety~\cite{li2015greenery,jiang2018brokenwindows,portnov2020lighting},
but does not by itself establish which local visual changes validly
shift perception in a given scene.

Following Jacobs and defensible-space theory~\cite{jacobs1961,newman1972},
we study a restricted version of this problem: whether a scene admits
multiple distinct \emph{single-lever} interventions that, evaluated
independently against the same unedited original, plausibly shift
perception.

We propose an \emph{interventional attribution} protocol that makes this
question testable. Rather than treating arbitrary image regions as the
explanatory unit, we define structured lever interventions drawn from an
editor-feasible vocabulary, each specifying a semantic concept, spatial
support, intervention direction, and constrained edit template.
Candidate edits are generated through prompt-conditioned diffusion
editing and retained only if a vision--language critic judges them to
preserve the same place, remain localised, and appear realistic and
plausible. Because generative editors may introduce correlated
non-target changes, each edit is treated as a hypothesis to be audited,
not as faithful evidence by default.

This paper does not claim validated causal effects on human perception; it introduces an auditable framework and pilot benchmark using validity-gated generation and an auxiliary perception proxy to prioritise interventions for future human evaluation.

Our empirical questions are: \emph{(i)} which lever families show directional improvement on an auxiliary perception proxy after validity auditing, and \emph{(ii)} across scenes, how many distinct single-lever interventions remain promising after screening and thresholding?

Our contributions:
\begin{enumerate}
    \item A \textbf{lever-based interventional attribution} framework
    for street-view perception, with the structured lever intervention
    $l=(c,s,d,\tau)$ as the explanatory atom.

    \item A \textbf{scene-specific candidate construction and evaluation
    pipeline} that instantiates levers from the intervention
    vocabulary, realises them through constrained prompt-only edits, and
    audits them for same-place preservation, locality, realism, and
    plausibility.

    \item A \textbf{reproducible pilot benchmark} across 50 scenes from five cities, characterising pipeline behaviour under prompt-only editing and reporting preliminary auxiliary patterns across lever families.
\end{enumerate}